% 附录 B：运行环境与可复现性
% ---------------------------------------------------------------------------
% 写作阶段填写。写明运行环境、依赖版本、随机种子与重跑方式，保证结论可复现。
% ---------------------------------------------------------------------------
%
% 建议结构（示例，按本题实际情况调整）：
%
% \begin{table}[htbp]
%   \centering
%   \begin{tabular}{ll}
%     \toprule
%     项目 & 配置 \\
%     \midrule
%     操作系统 & Windows 11（x86\_64） \\
%     Python & 3.12（uv 管理依赖） \\
%     关键依赖 & statsmodels / scikit-learn / lifelines / matplotlib \\
%     随机种子 & 全部随机过程固定 seed=42 \\
%     \bottomrule
%   \end{tabular}
% \end{table}
%
% 重跑方式：在仓库根目录执行 \texttt{uv sync} 安装依赖，再按
% \texttt{code/} 下脚本顺序运行（\texttt{load\_data.py} 之后逐问求解）。
% 交叉验证、Bootstrap 与测量误差扰动均固定种子，结果可复现。
